% Compte rendu hebdomadaire — stage CYBEL
% Compilation : pdflatex compte_rendu_semaine_2026-06-26.tex
\documentclass[11pt,a4paper]{article}

\usepackage[T1]{fontenc}
\usepackage[utf8]{inputenc}
\usepackage[french]{babel}
\usepackage{geometry}
\usepackage{booktabs}
\usepackage{longtable}
\usepackage{hyperref}
\usepackage{enumitem}
\usepackage{xcolor}
\usepackage{graphicx}

\geometry{margin=2.5cm}
\hypersetup{colorlinks=true, linkcolor=blue!60!black, urlcolor=blue!60!black}

\title{\textbf{Compte rendu d'activité — semaine du 23 au 26 juin 2026}\\[0.4em]
\large Projet CYBEL — Plateforme de commande robot CIOT TY1251D}
\author{[Nom du stagiaire]\\HESTIM Engineering \& Business School}
\date{Encadrant : Dr. Sridath Tula\\Rédigé le 26 juin 2026}

\begin{document}
\maketitle
\tableofcontents
\newpage

%==============================================================================
\section{Synthèse de la semaine}
%==============================================================================

Cette semaine a été consacrée à la \textbf{validation terrain} de l'option hybride \emph{Sentrymove / Deployment Tool + kiosque CYBEL} sur la carte \textbf{laboV2}, avec migration du parcours guidé (12 arrêts), synchronisation automatique des POI, correction de dysfonctionnements bloquants sur tablette Termux, et documentation opérationnelle pour le contrôleur labo.

\textbf{Livrables principaux :}
\begin{itemize}[noitemsep]
  \item Backend test \texttt{CybelVisitorKioskTest} opérationnel (port~8001, \texttt{~/cybel-test}).
  \item Sync POI ROS $\rightarrow$ \texttt{points.json} au démarrage kiosque et visite ; suppression des POI fantômes (ancienne carte).
  \item Correctifs navigation visite guidée (détection d'arrivée par proximité lorsque \texttt{nav\_status} reste à 601).
  \item Guides : \texttt{SENTRYMOVE\_POI\_SYNC.md}, \texttt{GUIDE\_CONTROLEUR\_POI.md}, mise à jour \texttt{TERRAIN.md}.
  \item Branche Git \texttt{feature/hybrid-sentrymove-kiosk} — une quinzaine de commits sur la période.
\end{itemize}

\textbf{État fonctionnel au 26/06 :}
\begin{itemize}[noitemsep]
  \item \textbf{Navigation point par point} (grille destinations) : validée terrain.
  \item \textbf{Visite guidée} : partiellement validée ; correctif déployé, re-test complet des 12 arrêts en attente.
  \item \textbf{Comparaison A/B} coords (port 8000) vs POI (port 8001) : en cours.
\end{itemize}

%==============================================================================
\section{Travaux réalisés (détail)}
%==============================================================================

\subsection{Migration carte laboV2 et parcours guidé}

\begin{itemize}[noitemsep]
  \item Passage à la carte \textbf{laboV2} dans \texttt{data/lab\_tour.json} (\texttt{map\_name: laboV2}, 12 arrêts avec \texttt{target\_point}).
  \item Alignement des noms POI sur le format Deployment Tool (MAJUSCULES, tirets).
  \item Mise à jour des présentations (\texttt{knowledgeV2-labV2.json}) et configuration kiosque test (\texttt{kiosk\_config.poi.json}).
  \item Déploiement USB/ADB sur tablette (\texttt{172.16.0.132}).
\end{itemize}

\subsection{Synchronisation POI (option hybride)}

Implémentation et documentation de la chaîne :
\[
\text{Deployment Tool} \rightarrow \text{marker\_manager (ROS)} \rightarrow \texttt{points.json} \rightarrow \text{kiosque CYBEL}
\]

\begin{itemize}[noitemsep]
  \item Module \texttt{sdk/poi\_sync.py}, \texttt{sdk/marker\_utils.py}, script \texttt{sync\_poi\_from\_robot.py}.
  \item API \texttt{POST /api/navigation/sync} (PC et \texttt{cybel\_lite}).
  \item \textbf{Sync automatique} : \texttt{GET /api/reception/destinations} (ouverture kiosque) et \texttt{POST /api/tour/start} (démarrage visite).
  \item \textbf{Élagage} : remplacement du cache (plus de fusion) — les POI absents de la carte ROS courante sont retirés.
  \item Tests unitaires : \texttt{tests/unit/test\_poi\_sync.py} (4 tests passants).
\end{itemize}

\subsection{Robustesse navigation et prérequis robot}

\begin{itemize}[noitemsep]
  \item Lecture \texttt{hard\_estop} depuis ROS (plus de valeur figée).
  \item Assouplissement des prérequis en statut 605 hors borne de charge.
  \item Messages d'erreur et hints de récupération (\texttt{tour\_navigation.py}).
\end{itemize}

\subsection{Documentation et formation}

\begin{itemize}[noitemsep]
  \item Guide contrôleur POI pour technicien labo HESTIM.
  \item Procédure courte \texttt{SENTRYMOVE\_POI\_SYNC.md}.
  \item Mise à jour \texttt{CHANGELOG.md} (v0.3.2), \texttt{PARITY\_CHECKLIST.md}.
\end{itemize}

%==============================================================================
\section{Problèmes rencontrés, hypothèses et solutions}
%==============================================================================

\begin{longtable}{@{}p{0.22\textwidth}p{0.28\textwidth}p{0.22\textwidth}p{0.22\textwidth}@{}}
\toprule
\textbf{Symptôme} & \textbf{Hypothèse} & \textbf{Solution appliquée} & \textbf{Statut} \\
\midrule
\endhead

Robot ne bouge pas ; message relocalisation & E-stop physique actif ; \texttt{nav\_status} 600/605 & Lecture \texttt{hard\_estop} ROS ; recovery 605 hors charge & Résolu \\

POI ancienne carte affichés sur kiosque & Cache \texttt{points.json} fusionné sans élagage & \texttt{merge\_point\_dicts} remplace par marqueurs ROS ; sync auto à l'ouverture & Résolu \\

Erreur à l'ouverture du kiosque (HTTP 500) & Import \texttt{sdk/\_\_init\_\_.py} $\rightarrow$ \texttt{pydantic} absent sur Termux lite & Chargement modules SDK par stub (\texttt{\_load\_sdk\_module\_from\_file}) ; déploiement \texttt{poi\_names.py} & Résolu \\

Backend TEST injoignable après déploiement & Délai démarrage Termux ; redémarrage via \texttt{su} sans Python & RUN\_COMMAND Termux ; scripts \texttt{start\_cybel\_test.sh} & Résolu \\

Visite guidée : erreur statut 601 alors que le robot est arrivé & Le châssis ne publie pas 602/603 malgré déplacement effectif (objectif POI accepté côté moteur) & Détection d'arrivée par \textbf{proximité} ($\leq$ 0{,}45\,m, \texttt{evaluate\_navigation\_arrival}) dans \texttt{cybel\_lite.py} & Correctif déployé ; validation 12 arrêts à faire \\

Échec intermittent arrêt LG-10 & \texttt{nav\_status} bloqué en 601 ; planificateur ROS silencieux & Même piste que ci-dessus ; traces \texttt{/api/tour/trace} & En observation \\

Nouveau POI (charge) invisible sur kiosque & \texttt{kiosk\_visible=false} par défaut hors parcours & Ouvrir grille destinations (sync auto) ; pas de redémarrage app nécessaire & Comportement attendu documenté \\

\bottomrule
\end{longtable}

\subsection{Hypothèse de travail retenue (stratégie S3)}

L'\textbf{option hybride POI Sentrymove} est la stratégie la plus fiable pour la navigation visiteur : le kiosque CYBEL réutilise les mêmes destinations nommées que l'application constructeur (\texttt{/tag\_manager/navi}), évitant les écarts de calibration des coordonnées brutes (\texttt{/navi\_goal}, stratégie S1). La comparaison A/B sur deux backends (8000 vs 8001) permet de quantifier cet écart sur le terrain HESTIM.

%==============================================================================
\section{Architecture de déploiement (rappel)}
%==============================================================================

\begin{center}
\begin{tabular}{@{}llll@{}}
\toprule
\textbf{Instance} & \textbf{Répertoire Termux} & \textbf{Port} & \textbf{APK} \\
\midrule
Production (coords) & \texttt{\textasciitilde/cybel} & 8000 & CybelVisitorKiosk \\
Test POI & \texttt{\textasciitilde/cybel-test} & 8001 & CybelVisitorKioskTest \\
\bottomrule
\end{tabular}
\end{center}

Robot : \texttt{192.168.20.22} (rosbridge 9090). Tablette : \texttt{172.16.0.132}.

%==============================================================================
\section{Perspectives — étapes suivantes et durées estimées}
%==============================================================================

\begin{longtable}{@{}p{0.05\textwidth}p{0.42\textwidth}p{0.12\textwidth}p{0.33\textwidth}@{}}
\toprule
\textbf{\#} & \textbf{Tâche} & \textbf{Durée} & \textbf{Critère de succès} \\
\midrule
\endhead

1 & Validation visite guidée complète 12 arrêts laboV2 (traces + corrections résiduelles) & 2--3 j & 12/12 arrêts sans erreur ; \texttt{nav\_status} ou proximité OK \\
2 & Comparaison A/B documentée (coords vs POI) : temps départ, précision, taux échec & 2 j & Tableau + extrait traces dans rapport \\
3 & Consolidation déploiement (\texttt{deploy\_termux.py}, bundle SDK complet incl. \texttt{poi\_names.py}) & 1 j & Déploiement reproductible en une commande \\
4 & Stabilisation E-stop / relocalisation (procédure affichée kiosque) & 1--2 j & Message opérateur clair ; recovery auto testée \\
5 & Prototype accueil par détection de présence (\texttt{/detected\_people\_array}) & 2--3 j & Sortie veille + TTS à l'approche \\
6 & Reconnaissance faciale (phase 2, optionnelle) & 8--10 j & Spike caméra tablette + annuaire visiteur \\
7 & Rédaction rapport / chapitres 5--7 + diagrammes & 5--7 j & Version soumise encadrant \\
8 & PostgreSQL + persistance visites (si scope maintenu) & 5 j & POI et historique survivent au redémarrage \\

\bottomrule
\end{longtable}

\textbf{Total estimé jusqu'à fin juillet} (hors vacances éventuelles) : \textbf{4 à 5 semaines} de travail effectif, selon disponibilité robot et reprises terrain.

%==============================================================================
\section{Besoin d'outils d'assistance au développement (IA)}
%==============================================================================

\subsection{Constat}

Le projet CYBEL combine des domaines à forte charge cognitive :
\begin{itemize}[noitemsep]
  \item rétro-ingénierie ROS et APK constructeur (milliers de fichiers Java décompilés) ;
  \item déploiement embarqué Android 7.1 / Termux / WebView ;
  \item cycles court debug terrain (tablette $\leftrightarrow$ robot $\leftrightarrow$ PC) avec journaux distribués.
\end{itemize}

Sans assistance logicielle adaptée, une part significative du temps est absorbée par la navigation dans le code, la rédaction de scripts de déploiement répétitifs et le diagnostic d'erreurs environnementales (ex.\ \texttt{ModuleNotFoundError: pydantic} sur Termux, identifié et corrigé en une session grâce à l'analyse automatisée des logs).

\subsection{Apport constaté cette semaine (exemples concrets)}

\begin{itemize}[noitemsep]
  \item \textbf{Diagnostic accéléré} : analyse du traceback \texttt{cybel-test-uvicorn.log} $\rightarrow$ cause racine (import \texttt{sdk/\_\_init\_\_}) en minutes au lieu d'une demi-journée de tests manuels.
  \item \textbf{Implémentation cohérente} : alignement de \texttt{cybel\_lite.py} sur \texttt{real\_robot.py} pour \texttt{evaluate\_navigation\_arrival} — réutilisation des tests existants (\texttt{test\_phase6.py}, 20 tests passants).
  \item \textbf{Documentation opérationnelle} : rédaction synchronisée des guides contrôleur et procédures sync POI (format exploitables par le labo HESTIM).
  \item \textbf{Déploiement} : commandes ADB, RUN\_COMMAND Termux, vérifications \texttt{curl} — exécution et itération sans rupture de concentration.
\end{itemize}

\subsection{Demande}

Il est demandé la prise en charge d'un \textbf{abonnement professionnel} à un environnement de développement assisté par IA, par exemple :

\begin{itemize}[noitemsep]
  \item \textbf{Cursor Pro} (IDE + agents intégrés au dépôt CYBEL) — environ 20\,€/mois ;
  \item ou \textbf{Claude Pro / équivalent} (analyse de longs journaux, rédaction technique) — ordre de grandeur similaire.
\end{itemize}

\textbf{Justification :}
\begin{enumerate}[noitemsep]
  \item \textbf{Gain de temps mesurable} : cette semaine, plusieurs blocages terrain (sync POI, pydantic Termux, visite guidée 601) ont été levés en une session de travail grâce à l'itération assistée ; estimation conservative \textbf{30 à 40\,\%} de temps gagné sur debug + doc pour la durée restante du stage.
  \item \textbf{Qualité} : respect des conventions du dépôt, tests unitaires maintenus, documentation à jour — réduction du risque de régression avant soutenance.
  \item \textbf{Continuité} : le stage se poursuit en juillet avec validation complète visite, comparaison A/B et rédaction du rapport ; l'outil amortit son coût sur 2--3 mois restants ($\ll$ coût d'une journée de blocage terrain).
  \item \textbf{Alternative moins favorable} : développement seul sans assistant — ralentissement prévisible sur les tâches transverses (déploiement, reverse engineering, rédaction LaTeX).
\end{enumerate}

\textbf{Recommandation} : privilégier \textbf{Cursor Pro} pour l'intégration directe au workflow Git / Termux / Python du projet ; compléter si besoin par Claude pour la rédaction du rapport final.

%==============================================================================
\section{Annexes}
%==============================================================================

\subsection{Commandes utiles (terrain)}

\begin{verbatim}
python scripts/sync_poi_from_robot.py --host 192.168.20.22
adb forward tcp:18001 tcp:8001
curl http://127.0.0.1:18001/api/reception/destinations
curl -X POST "http://127.0.0.1:18001/api/tour/start?lang=fr"
pytest tests/unit/test_poi_sync.py -q
\end{verbatim}

\subsection{Références documentaires}

\begin{itemize}[noitemsep]
  \item \texttt{docs/SENTRYMOVE\_POI\_SYNC.md}
  \item \texttt{docs/labo/GUIDE\_CONTROLEUR\_POI.md}
  \item \texttt{docs/cybel-conception/06-plan-hybride-sentrymove-kiosk.md}
  \item \texttt{docs/TOUR\_NAVIGATION.md}
\end{itemize}

\subsection{POI laboV2 (12 arrêts)}

\texttt{PORTE-LABO}, \texttt{CNC ROUTEUR}, \texttt{LG-10}, \texttt{IMPRIMANTE 3D}, \texttt{POINT-MACHINE}, \texttt{THERMOFORMAGE}, \texttt{EXTRUSION-SOUFFLAGE}, \texttt{POSTE-MACHINE}, \texttt{POSTE-REMPLISSAGE-BOUCHONNAGE}, \texttt{POSTE-ETIQUETAGE}, \texttt{GAMME-CONTROLE-QUALITE}, \texttt{SÉRIGRAPHIE}.

\vfill
\begin{center}
\small\textit{Document généré dans le cadre du stage CYBEL — HESTIM — juin 2026.}
\end{center}

\end{document}
